\documentclass[11pt,a4paper]{article}
\usepackage[utf8]{inputenc}
\usepackage[T1]{fontenc}
\usepackage[spanish]{babel}
\usepackage{graphicx}
\usepackage{float}
\usepackage[margin=2.2cm]{geometry}
\usepackage{enumitem}
\usepackage[hidelinks]{hyperref}

\title{Guía de observación\\EDA 01: estructura, cobertura y comparabilidad}
\author{Proyecto AgroCebada}
\date{}

\newcommand{\edafigure}[3]{
  \subsection{#1}
  \begin{figure}[H]
    \centering
    \includegraphics[width=0.96\textwidth,height=0.68\textheight,keepaspectratio]{\detokenize{#2}}
    \caption{#1}
  \end{figure}
  \textbf{Preguntas guía}
  #3
}

\begin{document}
\maketitle
\tableofcontents
\newpage

\section{Propósito}
Esta guía acompaña las salidas del primer EDA. Su objetivo es orientar la lectura
de la estructura temporal, cobertura, calidad y comparabilidad de Sentinel-2,
Landsat y Planet. Las gráficas son descriptivas: no implican equivalencia entre
sensores ni justifican todavía imputación, escalamiento o selección de variables.

\section{Cobertura temporal y por parcela}

\edafigure{Fechas disponibles por sensor}
{../../../outputs/eda/01_tablas_maestras/01_eda_temporal/01_timeline_sensores.png}{
\begin{itemize}[nosep]
  \item ¿Qué sensores cubren todo el periodo y cuál está limitado a 2025?
  \item ¿Se observan huecos, temporadas densas o cambios de frecuencia?
  \item ¿Qué periodos permiten comparaciones temporales justas entre sensores?
\end{itemize}}

\edafigure{Intervalos entre observaciones consecutivas}
{../../../outputs/eda/01_tablas_maestras/01_eda_temporal/02_intervalos_observaciones.png}{
\begin{itemize}[nosep]
  \item ¿Cuál es el intervalo típico de revisita por sensor?
  \item ¿La distribución tiene colas largas que indiquen interrupciones?
  \item ¿La frecuencia observada es suficientemente regular para comparar trayectorias?
\end{itemize}}

\edafigure{Cobertura por parcela}
{../../../outputs/eda/01_tablas_maestras/01_eda_temporal/03_cobertura_parcelas.png}{
\begin{itemize}[nosep]
  \item ¿Cuántas observaciones suele tener una parcela en cada sensor?
  \item ¿Existen parcelas con cobertura muy inferior o superior al resto?
  \item ¿Las diferencias entre sensores se explican por su rango temporal?
\end{itemize}}

\edafigure{Heatmap de cobertura de Sentinel-2}
{../../../outputs/eda/01_tablas_maestras/01_eda_temporal/04_heatmap_sentinel_2.png}{
\begin{itemize}[nosep]
  \item ¿La cobertura mensual es uniforme entre parcelas?
  \item ¿Qué meses muestran huecos compartidos por muchas parcelas?
  \item ¿Los patrones ausentes parecen sistemáticos o aislados?
\end{itemize}}

\edafigure{Heatmap de cobertura de Landsat}
{../../../outputs/eda/01_tablas_maestras/01_eda_temporal/05_heatmap_landsat.png}{
\begin{itemize}[nosep]
  \item ¿Qué tan heterogénea es la cobertura entre parcelas?
  \item ¿Hay bloques temporales sin observaciones?
  \item ¿Qué parcelas requieren atención antes de un análisis longitudinal?
\end{itemize}}

\edafigure{Heatmap de cobertura de Planet}
{../../../outputs/eda/01_tablas_maestras/01_eda_temporal/06_heatmap_planet.png}{
\begin{itemize}[nosep]
  \item ¿La cobertura durante 2025 es continua para las parcelas mostradas?
  \item ¿Qué meses presentan menor disponibilidad?
  \item ¿La cobertura de Planet coincide con meses comparables de los otros sensores?
\end{itemize}}

\section{Calidad y nubosidad}

\edafigure{Distribución de nubosidad}
{../../../outputs/eda/01_tablas_maestras/01_eda_temporal/07_distribucion_nubosidad.png}{
\begin{itemize}[nosep]
  \item ¿La nubosidad se concentra en 0, 100 o valores intermedios?
  \item ¿Las distribuciones tienen formas comparables entre sensores?
  \item ¿Los valores extremos pueden corresponder a reglas de procesamiento distintas?
\end{itemize}}

\edafigure{Nubosidad por mes}
{../../../outputs/eda/01_tablas_maestras/01_eda_temporal/08_nubosidad_por_mes.png}{
\begin{itemize}[nosep]
  \item ¿Existe estacionalidad en la nubosidad?
  \item ¿Qué meses presentan mayor mediana o dispersión?
  \item ¿La comparación está afectada por los distintos años disponibles por sensor?
\end{itemize}}

\edafigure{Cobertura mensual frente a nubosidad}
{../../../outputs/eda/01_tablas_maestras/01_eda_temporal/09_cobertura_vs_nubosidad.png}{
\begin{itemize}[nosep]
  \item ¿Los meses con mayor nubosidad observada tienen menos registros?
  \item ¿La relación cambia según el sensor?
  \item ¿Puede hablarse solo de asociación y no de causalidad, dado que no se conoce la nubosidad de capturas ausentes?
\end{itemize}}

\section{Trayectorias temporales representativas}

\edafigure{Trayectorias de NDVI promedio}
{../../../outputs/eda/01_tablas_maestras/01_eda_temporal/10_trayectoria_ndvi_promedio.png}{
\begin{itemize}[nosep]
  \item ¿Se reconocen ciclos de crecimiento y senescencia?
  \item ¿Los sensores reproducen tendencias similares en periodos coincidentes?
  \item ¿Hay saltos abruptos que ameriten revisar nubosidad o calidad?
\end{itemize}}

\edafigure{Trayectorias de LAI promedio}
{../../../outputs/eda/01_tablas_maestras/01_eda_temporal/11_trayectoria_lai_promedio.png}{
\begin{itemize}[nosep]
  \item ¿Los máximos de LAI ocurren en fechas semejantes a los máximos de NDVI?
  \item ¿Existe un desplazamiento sistemático de escala entre sensores?
  \item ¿Hay valores o cambios biológicamente poco plausibles?
\end{itemize}}

\edafigure{Trayectorias de NDWI promedio}
{../../../outputs/eda/01_tablas_maestras/01_eda_temporal/12_trayectoria_ndwi_promedio.png}{
\begin{itemize}[nosep]
  \item ¿Cómo cambia la señal de humedad a lo largo del ciclo?
  \item ¿Los cambios coinciden temporalmente con NDVI y LAI?
  \item ¿Qué parcelas muestran comportamientos atípicos?
\end{itemize}}

\edafigure{Trayectorias de ENDWI promedio}
{../../../outputs/eda/01_tablas_maestras/01_eda_temporal/13_trayectoria_endwi_promedio.png}{
\begin{itemize}[nosep]
  \item ¿ENDWI confirma los patrones observados en NDWI?
  \item ¿Presenta mayor ruido o sensibilidad a valores extremos?
  \item ¿Hay periodos sin datos que impidan interpretar la trayectoria?
\end{itemize}}

\edafigure{Trayectorias de MSI promedio}
{../../../outputs/eda/01_tablas_maestras/01_eda_temporal/14_trayectoria_msi_promedio.png}{
\begin{itemize}[nosep]
  \item ¿Los aumentos de estrés hídrico coinciden con descensos de NDVI?
  \item ¿La dirección y magnitud de los cambios son consistentes entre parcelas?
  \item ¿Hay observaciones aisladas que dominen la interpretación?
\end{itemize}}

\edafigure{Trayectorias de EMSI promedio}
{../../../outputs/eda/01_tablas_maestras/01_eda_temporal/15_trayectoria_emsi_promedio.png}{
\begin{itemize}[nosep]
  \item ¿EMSI aporta un patrón diferente o redundante respecto a MSI?
  \item ¿Los picos coinciden con periodos de baja humedad?
  \item ¿La trayectoria es estable frente a posibles observaciones contaminadas?
\end{itemize}}

\edafigure{Trayectorias de NDDI promedio}
{../../../outputs/eda/01_tablas_maestras/01_eda_temporal/16_trayectoria_nddi_promedio.png}{
\begin{itemize}[nosep]
  \item ¿Los máximos sugieren periodos coherentes de sequedad?
  \item ¿NDDI se mueve de manera opuesta a los índices de humedad?
  \item ¿Los cambios abruptos coinciden con alta nubosidad o huecos temporales?
\end{itemize}}

\edafigure{Trayectorias de ENDDI promedio}
{../../../outputs/eda/01_tablas_maestras/01_eda_temporal/17_trayectoria_enddi_promedio.png}{
\begin{itemize}[nosep]
  \item ¿ENDDI confirma los periodos secos detectados con NDDI?
  \item ¿Qué diferencias se observan entre las parcelas seleccionadas?
  \item ¿La disponibilidad temporal permite comparar ciclos completos?
\end{itemize}}

\section{Comparación preliminar Planet--Sentinel-2}

\edafigure{Comparación exacta de NDVI}
{../../../outputs/eda/01_tablas_maestras/01_eda_temporal/18_planet_sentinel_ndvi_promedio.png}{
\begin{itemize}[nosep]
  \item ¿Los puntos siguen una relación aproximadamente lineal?
  \item ¿Existe sesgo sistemático entre Planet y Sentinel-2?
  \item ¿La dispersión cambia para valores bajos y altos de NDVI?
\end{itemize}}

\edafigure{Comparación exacta de EVI}
{../../../outputs/eda/01_tablas_maestras/01_eda_temporal/19_planet_sentinel_evi_promedio.png}{
\begin{itemize}[nosep]
  \item ¿EVI presenta mayor o menor concordancia visual que NDVI?
  \item ¿Hay grupos, valores extremos o diferencias de escala?
  \item ¿Cuánto de la dispersión podría relacionarse con resolución o procesamiento?
\end{itemize}}

\edafigure{Comparación exacta de LAI}
{../../../outputs/eda/01_tablas_maestras/01_eda_temporal/20_planet_sentinel_lai_promedio.png}{
\begin{itemize}[nosep]
  \item ¿Se observa un desplazamiento consistente entre ambos sensores?
  \item ¿La relación conserva el orden aun si las escalas son distintas?
  \item ¿Qué evidencia adicional sería necesaria antes de armonizar LAI?
\end{itemize}}

\section{Respuestas y hallazgos}

\subsection{Cobertura temporal y por parcela}
\begin{enumerate}
  \item \textbf{Timeline.} Sentinel-2 y Landsat cubren de enero de 2022 a
  diciembre de 2025; Planet sólo cubre 2025. Por ello, la comparación de los
  tres sensores debe limitarse a 2025. Hay 301 fechas únicas de Sentinel-2, 392
  de Landsat y 311 de Planet. Visualmente no hay interrupciones anuales grandes,
  aunque la densidad de fechas es distinta.

  \item \textbf{Intervalos.} La revisita mediana observada es de 5 días para
  Sentinel-2, 7 días para Landsat y 1 día para Planet. El percentil 99 es de 10,
  8 y 5 días, respectivamente; los máximos son 10, 16 y 19 días. Planet es el
  más frecuente, pero presenta algunos huecos aislados más largos. Las series
  no deben compararse por posición de fila, sino por fecha o mediante ventanas
  temporales explícitas.

  \item \textbf{Cobertura por parcela.} Sentinel-2 es muy uniforme, con
  298--301 observaciones por parcela. Planet también es relativamente uniforme,
  con 233--251. Landsat es heterogéneo: 195--391, mediana de 211 y media de
  246.8. La diferencia no se explica por el rango temporal, pues todas las
  parcelas Landsat comparten fechas inicial y final; refleja disponibilidad
  desigual dentro del periodo.

  \item \textbf{Heatmaps.} Los tres sensores alcanzan las 197 parcelas en cada
  mes representado. Sentinel-2 muestra bloques mensuales muy regulares y Planet
  una cobertura densa durante 2025. Landsat presenta la mayor desigualdad en el
  número de observaciones por celda parcela--mes. No hay meses completos
  ausentes para todas las parcelas, pero sí diferencias internas que deben
  conservarse al analizar trayectorias.
\end{enumerate}

\subsection{Calidad y nubosidad}
\begin{enumerate}[resume]
  \item \textbf{Distribución.} La variable está fuertemente concentrada en los
  extremos. El valor 0 representa 62.24\% de Sentinel-2, 59.42\% de Landsat y
  63.42\% de Planet; el valor 100 representa 23.99\%, 31.87\% y 25.22\%,
  respectivamente, entre registros válidos. Las medianas son 0 en los tres
  sensores. Esta forma sugiere una variable derivada de reglas o máscaras y no
  una medición continua plenamente comparable.

  \item \textbf{Estacionalidad.} En 2025, junio tiene la mayor nubosidad media
  en Sentinel-2 (66.0\%) y Landsat (74.4\%); Planet alcanza su máximo en
  septiembre (56.8\%). Los mínimos ocurren en noviembre para Sentinel-2
  (1.8\%), abril para Landsat (0.3\%) y marzo para Planet (6.5\%). La gráfica
  con todos los años mezcla periodos distintos, por lo que la comparación justa
  entre los tres debe apoyarse en la versión 2025 del EDA 02.

  \item \textbf{Cobertura frente a nubosidad.} No aparece una relación negativa
  uniforme. La correlación mensual entre cantidad de filas y nubosidad mediana
  es aproximadamente $-0.05$ en Sentinel-2, $0.15$ en Landsat y $-0.44$ en
  Planet. Es una asociación descriptiva y no prueba que las nubes causen
  ausencias: la nubosidad de una captura no registrada no está disponible.
\end{enumerate}

\subsection{Trayectorias representativas}
Las parcelas seleccionadas automáticamente son AGC\_099, AGC\_154 y AGC\_012,
correspondientes aproximadamente a coberturas total baja, mediana y alta.
\begin{enumerate}[resume]
  \item \textbf{NDVI.} Las series muestran ciclos de valores bajos y picos de
  vegetación, pero también descensos y recuperaciones abruptas. En las parcelas
  elegidas, Sentinel-2 va de $-0.020$ a $0.766$ y Planet de $0.116$ a $0.852$.
  Durante 2025 las tendencias coinciden en términos generales, aunque Planet
  suele quedar por encima; los saltos requieren revisar nubosidad y separación
  entre fechas.

  \item \textbf{LAI.} Los picos acompañan de forma general los periodos de NDVI
  alto, pero las escalas no coinciden: Sentinel-2 va de $-0.180$ a $1.806$ y
  Planet de $0.604$ a $3.267$. El LAI negativo de Sentinel-2 es poco plausible
  físicamente y el desplazamiento entre sensores impide tratarlos como una sola
  escala sin calibración.

  \item \textbf{NDWI y ENDWI.} Sólo están disponibles en Sentinel-2 dentro de
  estas tablas. NDWI varía de $-0.295$ a $0.676$ y ENDWI de $-0.255$ a $0.718$
  en las parcelas elegidas. Ambos describen patrones parecidos de humedad, pero
  ENDWI no elimina los cambios bruscos ni resuelve los huecos producidos por
  valores faltantes.

  \item \textbf{MSI y EMSI.} También son exclusivos de Sentinel-2. Sus rangos
  son 0.193--1.846 y 0.164--1.697. Los niveles altos tienden a aparecer cuando
  los índices de vegetación/humedad son bajos, como se espera para estrés
  hídrico; aun así, los picos aislados deben contrastarse con nubosidad.

  \item \textbf{NDDI y ENDDI.} Presentan rangos más estrechos, de
  $-0.232$ a $0.238$ y de $-0.237$ a $0.264$. Ambos resumen un patrón semejante
  de sequedad y son en buena medida redundantes visualmente en estas parcelas.
  Antes de interpretarlos agronómicamente se debe revisar si sus extremos
  coinciden con capturas nubladas o con discontinuidades temporales.

  \item \textbf{Disponibilidad de índices.} En las tres parcelas hay 227
  faltantes para cada trayectoria Sentinel-2 mostrada y 182 en las variables
  Planet. Por ello, las líneas conectan puntos válidos separados y no implican
  mediciones continuas entre fechas.
\end{enumerate}

\subsection{Comparación Planet--Sentinel-2}
Las comparaciones con valores válidos contienen 8,104 pares exactos parcela--fecha.
\begin{enumerate}[resume]
  \item \textbf{NDVI.} La asociación es alta ($r=0.974$), pero Planet supera en
  promedio a Sentinel-2 por 0.083 unidades; el MAE es 0.083. La diferencia es
  sistemática y no sólo ruido aleatorio.

  \item \textbf{EVI.} Presenta $r=0.958$, diferencia media Planet menos
  Sentinel-2 de $-0.0013$ y MAE de 0.0357. Es la variable con menor sesgo medio,
  aunque todavía existe dispersión y no se demuestra equivalencia.

  \item \textbf{LAI.} Es la comparación menos concordante de las tres
  ($r=0.910$) y muestra el mayor desplazamiento: Planet supera a Sentinel-2 en
  1.062 unidades en promedio, igual prácticamente al MAE. Esto evidencia un
  sesgo de escala que exige calibración antes de combinar sensores.
\end{enumerate}

\section{Cierre sugerido}
Al terminar la revisión, documentar: periodos comparables, parcelas con cobertura
insuficiente, meses problemáticos, variables sensibles a nubosidad, diferencias
sistemáticas entre sensores y anomalías que requieran validación. No concluir que
dos sensores son intercambiables únicamente por una correlación visual alta.

\end{document}

